
Automated research pipelines with feasibility constraints fail when constraint violations propagate undetected through phase boundaries to expensive downstream implementation phases. Existing schema validation checks structural properties (types, required fields) but cannot express semantic constraints (keyword patterns) or compositional rules (cross-field logic, state-based membership). This work introduces a three-layer contract-based validation architecture (schema + pattern + compositional contracts) applied at research pipeline phase boundaries. The approach is validated on placeholder hypothesis content with typed interfaces under four feasibility constraints: no synthetic data generation, no human evaluation, standard datasets only, and existing benchmarks only. On a 100-case placeholder corpus (80 valid, 20 with embedded violations), contract-based validation achieved 100\% downstream failure reduction versus schema-only baseline (20/20 violations caught at boundaries versus 0/20). Contracts force explicit checking of constraints schema validation cannot express (100 percentage point expressiveness gap: 4/4 constraints versus 0/4), multi-layer detection achieves 88\% violation detection at boundaries versus 40\% for schema-only (48 percentage point gap), and early boundary validation prevents all constraint violations from propagating to Phase 4/5 implementation. The contract layer requires 25-27 lines of code with negligible execution overhead (0.01ms). External validity to substantive research content remains unproven; results demonstrate that constraint enforcement logic functions correctly on typed interfaces but do not establish generalization to production workflows with complex semantic dependencies.
